% Consolidated guide-selftest01 from f8f4709b0ffac9da4c1fecbf2ce873002f6da6cc; source fragments preserved.

% === BEGIN guide-selftest01-choices.tex blob 394051b8be498b943b41a287443388314c6ba5b2 ===
% 自测题1：基础力学选择1.1.1—1.1.10
% 《大学物理学学习指导》第5版，印刷45—46页／PDF53—54页。
\subsection{自测题1：运动学、动量和弹簧选择题}

\sourcetest{1.1.1}{45}{53}
\problemtext 沿$x$轴运动的质点有$x(t)=3t-5t^3+6$（SI），判断加速度方向和变化性质。
\workings
速度$v_x=3-15t^2$，加速度
\[
\boxed{a_x=-30t}.
\]
对$t>0$它指向$x$负方向且大小随时间增大，并非匀加速，故选\textbf{D}。
\crosscheck $t=0$瞬间$a_x=0$；说“向负方向”指题意通常考察的$t>0$区间。

\sourcetest{1.1.2}{45}{53}
\problemtext 曲线运动中用位矢$\bm r$、路程$s$、速度大小$v$及切向加速度大小$a_\tau$判断四个微分式。
\workings
速度的定义为$\bm v=d\bm r/dt$，而标量路程的变化率
$ds/dt=|\bm v|=v$；加速度向量$d\bm v/dt=\bm a$，
但其模通常并非切向加速度大小：
\[
|\bm a|=\sqrt{a_\tau^2+a_n^2}\ne a_\tau.
\]
因此只把$ds/dt=v$视为题面无歧义的标量恒等式，
对应\textbf{D（仅第3项）}。
\crosscheck 原书其余选项将速度矢量、速度标量与曲线运动的切向量混用。
若将第2式明确定义为\emph{矢量}等式$d\bm r/dt=\bm v$，它当然成立；
但此与原题印刷未加粗的$v$记号不同，属于须保留的记号歧义。

\sourcetest{1.1.3}{45}{53}
\problemtext 速度满足$dv/dt=-kv^2t$，$k>0$，初值$v(0)=v_0>0$，求$v(t)$。
\workings
分离变量
$dv/v^2=-kt\,dt$，
积分得
\[
-\frac1v=-\frac{kt^2}{2}-\frac1{v_0},
\qquad
\boxed{\frac1v=\frac{kt^2}{2}+\frac1{v_0}}.
\]
选\textbf{C}。
\crosscheck $t=0$恢复初值；对$t>0$、$v>0$时速度下降。

\sourcetest{1.1.4}{45}{53}
\problemtext 水平粗糙面上的物块受固定大小$F$、向右上方夹角$\theta$的拉力，摩擦系数$\mu$，求最大水平加速度时的$\theta$。
\workings
在物块始终与地面接触的范围内
$N=mg-F\sin\theta\ge0$，
\[
ma_x=F\cos\theta-\mu(mg-F\sin\theta).
\]
对角度求导，
$-F\sin\theta+\mu F\cos\theta=0$，
给
\[
\boxed{\tan\theta=\mu}.
\]
按题面答案选\textbf{C}。
\crosscheck 必须另检验所求角对应$N\ge0$；
若$F$很大使物体离地，摩擦模型不再适用。

\sourcetest{1.1.5}{45}{53}
\problemtext 光滑半球形碗内半径$R=0.10$m，
小球相对碗静止于离碗底$h=0.04$m处，
碗绕竖直轴匀速旋转，求角速度（取$g=9.8\,\mathrm{m/s^2}$）。
\workings
球心比小球高$R-h=0.06$m，
转动半径为$\rho=\sqrt{2Rh-h^2}$。
法向支持力沿球面法线，其竖直平衡与水平向心分量给
\[
N\frac{R-h}{R}=mg,\qquad
N\frac{\rho}{R}=m\omega^2\rho.
\]
消去$N$后
\[
\boxed{\omega=\sqrt{\frac g{R-h}}
=\sqrt{\frac{9.8}{0.06}}
\simeq12.78\,\mathrm{rad/s}}.
\]
选\textbf{A（约13 rad/s）}。
\crosscheck 原书位置高度是离\emph{碗底}4cm，不是离球心4cm。

\sourcetest{1.1.6}{45--46}{53--54}
\problemtext 质量2kg粒子在$\bm F=12t\bm i$（SI）作用下从原点静止开始运动，求$t=3$s时动量。
\workings
动量定理$\bm p(3)=\int_0^3\bm F\,dt$，
\[
\boxed{\bm p=12\int_0^3t\,dt\,\bm i
=54\bm i\,\mathrm{kg\,m/s}}.
\]
选\textbf{B}。
\crosscheck 力对时间积分直接为动量，不需再次乘物体质量。

\sourcetest{1.1.7}{46}{54}
\problemtext 小球在水平面内半径$R$的圆周作逆时针匀速率$v$运动，
从右侧点A（速度向上）走半周至左侧点B（速度向下），求动量增量。
\workings
图中$\bm p_A=mv\bm j$，$\bm p_B=-mv\bm j$，因此
\[
\boxed{\Delta\bm p=\bm p_B-\bm p_A=-2mv\bm j}.
\]
选\textbf{B}。
\crosscheck 动量是矢量；走半周后速率没变，但动量发生大小$2mv$的变化。

\sourcetest{1.1.8}{46}{54}
\problemtext 两个弹簧$k_A,k_B$串联，静止悬挂物块；求两弹簧弹性势能之比。
\workings
轻弹簧串联受相同拉力$T$，伸长量$x_i=T/k_i$，
\[
U_i=\frac12 k_ix_i^2=\frac{T^2}{2k_i}.
\]
故
\[
\boxed{\frac{U_A}{U_B}=\frac{k_B}{k_A}}.
\]
选\textbf{C}。
\crosscheck 同一拉力下较软弹簧储存的能量更大，不能用$k_A/k_B$直接代替。

\sourcetest{1.1.9}{46}{54}
\problemtext 光滑水平面上等质量物体P冲向带弹簧挡板的静止物体Q，碰后P停下、Q以原来P速度离开。何时弹簧压缩量最大？
\workings
压缩到最大时弹簧伸缩速率为零，
即两物体瞬时\emph{速度相等}，
相对动能全部转换成弹簧弹性势能：
\[
\boxed{v_P=v_Q}.
\]
选\textbf{B}。
\crosscheck P最终停下是碰撞完成后的事件，
不能把它等同最大压缩的时刻。

\sourcetest{1.1.10}{46}{54}
\problemtext 质量$M$重物悬于定滑轮右端，质量$m=M/2$的人抓左端绳向上爬，其相对绳的向上加速度为$a_0$，求人对地加速度（向上正）。
\workings
设绳左段加速度为$a_L$，重物向上的加速度为$-a_L$，
张力为$T$。对人和重物：
\[
ma_P=T-mg,\qquad
M(-a_L)=T-Mg,\qquad a_P-a_L=a_0.
\]
联立得
\[
a_P=\frac{M a_0+(M-m)g}{M+m}.
\]
代入$m=M/2$，
\[
\boxed{a_P=\frac{2a_0+g}{3}}.
\]
选\textbf{A}。
\crosscheck 必须明确“相对绳上爬”与“对地加速度”不同；
不能直接写$a_P=a_0$。
% === END guide-selftest01-choices.tex ===

% === BEGIN guide-selftest01-blanks.tex blob 41049fb796fca5f2f314edbd6d054e9f1a849366 ===
% 《大学物理学学习指导（第5版）》自测题1填空1.2.1—1.2.8
% 原书印刷46—47／扫描PDF54—55页。
\subsection{自测题1：相对运动、冲量及能量填空}

\sourcetest{1.2.1}{46--47}{54--55}
\problemtext 两辆车同时从原点出发，沿同一正方向运动，
$ x_A=4t+t^2$，$x_B=2t^2+2t^3$（SI）。
求刚出发时谁在前、离开后何时再次同地、何时两车瞬时速度相等。
\workings
$t\to0^+$时$x_A\sim4t>x_B\sim2t^2$，因此\textbf{A车在前}。
再次同地$x_A=x_B$，
$t(4-t-2t^2)=0$，正根为
\[
\boxed{t_{\rm meet}=\frac{\sqrt{33}-1}{4}\,\mathrm s
\simeq1.186\,\mathrm s}.
\]
两车速度$v_A=4+2t$、$v_B=4t+6t^2$，令其相等：
$3t^2+t-2=0$，正根
\[
\boxed{t_{v_A=v_B}=\frac23\,\mathrm s}.
\]
\crosscheck 两车位置相等和速度相等是两个不同事件，
不能用同一个时间填这两空。

\sourcetest{1.2.2}{47}{55}
\problemtext 列车沿水平以10m/s行驶，地面看雨竖直下落，
车窗雨迹偏离竖直方向30°，求雨相对地面和列车的速率。
\workings
设雨相对地面速率为$u$，相对车厢
$\bm v_{\rm rel}=(-10,-u)$，
按与竖直夹角
$\tan30^\circ=10/u$，故
\[
\boxed{u=10\sqrt3\simeq17.32\,\mathrm{m/s}},
\qquad
\boxed{|\bm v_{\rm rel}|=\sqrt{10^2+u^2}
=20.0\,\mathrm{m/s}}.
\]
\crosscheck 夹角是与\emph{竖直}而不是水平所成30°；
雨迹的斜向由车速方向确定。

\sourcetest{1.2.3}{47}{55}
\problemtext 静止小球B质量$m$，被水平绳AB和与竖直夹角$\theta$的斜绳BC平衡悬挂。
若瞬间剪断水平绳AB，求剪断瞬间斜绳BC的张力。
\workings
原来小球静止，剪断AB瞬间速度仍为零；
斜绳长度约束下沿BC方向的向心加速度$v^2/L=0$。
把重力投影到斜绳方向：
\[
\boxed{T=mg\cos\theta}.
\]
\crosscheck 剪断后有切向加速度$g\sin\theta$，
不能因为初速为零就误判所有方向加速度均为零。

\sourcetest{1.2.4}{47}{55}
\problemtext 30kg物块向东10m/s，
20kg物块向北20m/s，二者完全非弹性碰合，
求共同速率与方位。
\workings
总动量分量为$(p_E,p_N)=(300,400)$kg\,m/s，
总质量50kg，故
\[
\boxed{\bm v=(6\,\bm e_E+8\,\bm e_N)\,\mathrm{m/s}},
\qquad \boxed{v=10.0\,\mathrm{m/s}}.
\]
方向$\tan\phi=8/6=4/3$，
$\boxed{\phi\simeq53.13^\circ}$，即\textbf{东偏北53.13°}。
\crosscheck 动量矢量守恒，机械动能在完全非弹性碰撞中不守恒。

\sourcetest{1.2.5}{47}{55}
\problemtext 圆锥摆小球在水平面匀速转动整一周，
求动量增量、重力冲量与绳张力冲量大小。
\workings
小球回到同样速度方向，故
\[
\boxed{|\Delta\bm p|=0}.
\]
周期$T=2\pi/\omega$，
重力冲量$\bm J_g=-mgT\bm e_z$。
由总冲量为零，
$\bm J_T=-\bm J_g=+mgT\bm e_z$；
所以
\[
\boxed{|\bm J_g|=mgT},\qquad
\boxed{|\bm J_T|=mgT}.
\]
\crosscheck 虽然瞬时绳张力有水平向心分量，
整周的水平冲量积分为零，不可把张力大小直接乘周期。

\sourcetest{1.2.6}{47}{55}
\problemtext 光滑水平面上质量$m$物体在恒力$F$下从静止运动时间$t$，
求$F$的功；另一惯性观察者沿与$F$相反方向以恒速$v_0$运动，
问他测量的同一段时间内力做的功。
\workings
地面系位移$x=Ft^2/(2m)$，
\[
\boxed{W=\frac{F^2t^2}{2m}}.
\]
观察者的速度为$-v_0$，故物体相对该观察者
$x'=x+v_0t$（沿$F$正方向）。
于是
\[
\boxed{W'=\frac{F^2t^2}{2m}+Fv_0t}.
\]
\crosscheck 若把“观察者以$v_0$沿与$F$相反方向运动”改为
“物体起初相对观察者朝相反方向运动”，符号会改变。
功依参考系，不是伽利略不变量；此处严格采用原题观察者方向。

\sourcetest{1.2.7}{47}{55}
\problemtext 冰块从静止沿与水平成30°的无摩擦屋顶滑下10m后离开屋檐，
屋檐离地10m，取$g=10$，求离檐到落地的水平飞行距离。
\workings
沿斜面重力做功，离檐速率
$v=\sqrt{2g(10\sin30^\circ)}=10$ m/s。
离檐水平、竖直初速度分别为
$v_x=10\cos30^\circ=5\sqrt3$，
$v_y=-10\sin30^\circ=-5$。
从屋檐高度10m到地面的运动方程
$0=10-5t-5t^2$给$t=1$ s。
故
\[
\boxed{x=v_xt=5\sqrt3\simeq8.66\,\mathrm m}.
\]
\crosscheck 题目仅问\emph{脱离屋檐后}的水平位移，
不包括沿屋顶滑行10m时的水平投影。

\sourcetest{1.2.8}{47}{55}
\problemtext 两质点系统只受质点间相互引力且外力矢量和为零。
判断动量与机械能是否必然守恒。
\workings
内力按牛顿第三定律成对相消，
外力合力为零，
所以系统总动量满足
\[
\boxed{d\bm P/dt=0}.
\]
若系统只受互相引力这种保守内力，
系统机械能亦守恒。
但题面的“只有万有引力”严格假设确实确保保守，
故
\[
\boxed{\bm P=\mathrm{const}},\qquad
\boxed{E_{\rm mech}=\mathrm{const}}.
\]
应选原书\textbf{A（动量与机械能均守恒）}。
\crosscheck 如果把题设放宽为\emph{任意内力}，
即使外力合力为零，内摩擦仍会破坏机械能守恒。
本题明确是\emph{只有万有引力}，因此不可用一般内非保守力
错误地排除机械能守恒。
% === END guide-selftest01-blanks.tex ===

% === BEGIN guide-selftest01-solutions.tex blob 193663218320ba2c5b148274855ada0cfdb1dfe0 ===
% 《大学物理学学习指导》第5版自测题1：计算1.3.1—1.3.5，改错1.4.1。
% 原书印刷47—48／PDF物理55—56页；图示按原书方向转述。
\subsection{自测题1：动量、功与机械能计算}

\sourcetest{1.3.1}{47}{55}
\problemtext 质量$m$的小块从高度$h$的光滑曲面滑下，
曲面固定在质量$M$的可水平自由滑动楔块上。
楔面从左侧$x=0$延伸至右端$x=L$，
初始系统静止，求滑块落至底部期间楔块水平移动距离。
\workings
取楔块相对地面位移$\Delta X$（向右正），
小块的绝对水平位移是$\Delta X+L$。
全程没有外界水平力，系统质心的水平位置不变，
\[
M\Delta X+m(\Delta X+L)=0.
\]
因此
\[
\boxed{\Delta X=-\frac{m}{M+m}L}.
\]
楔块向左移动，距离大小为$mL/(M+m)$。
\crosscheck 位移由水平质心守恒给出，
\emph{与}曲面下落高度$h$的具体数值无关；
$h$决定速度与能量，不决定上述质心位移结果。

\sourcetest{1.3.2}{47}{55}
\problemtext 两艘质量同为$M$的船初始在静水中静止，
质量$m$的人以对地水平速度$v$从第一艘跳到第二艘，
忽略水阻，求两艘船在跳离及落入后各自速度。
\workings
以跳跃方向为正。
人刚离开船1时，船1与人的总水平动量守恒：
$MV_1+mv=0$，故
\[
\boxed{V_1=-\frac mM v}.
\]
人登上最初静止的船2并与它共同运动时，
对人+船2的瞬时完全非弹性相互作用应用动量守恒，
$(M+m)V_2=mv$，故
\[
\boxed{V_2=\frac{m}{M+m}v}.
\]
\crosscheck 船1向与人相反方向退行，船2沿人跳的方向前进。
人给定的$v$是\emph{相对于地面}，不能误当相对船的速度。

\sourcetest{1.3.3}{47--48}{55--56}
\problemtext 皮带以恒速$v$沿水平方向运行，
无初始水平速度的砂从静止漏斗落到带上。
单位时间落砂质量$\dot m=dM/dt$，
求维持皮带速度所需水平拉力和电机输出机械功率；
再取$\dot m=20\,\mathrm{kg/s},v=1.5\,\mathrm{m/s}$。
\workings
在$dt$内有$dm=\dot m\,dt$的砂从水平速度0
加速到皮带速度$v$，动量增量为$v\,dm$，
所以皮带必须施于砂的平均水平力
$F=v\,dm/dt=\dot m v$，
相应电机对皮带的机械功率
\[
\boxed{P=Fv=\dot m v^2}.
\]
数值为
\[
\boxed{F=20(1.5)=30\,\mathrm N},\qquad
\boxed{P=20(1.5)^2=45\,\mathrm W}.
\]
\crosscheck 砂获得的平动动能率仅为
$\tfrac12\dot m v^2=22.5$W；
其余机械功率在砂与带相对滑动时转为热。
因此\emph{电机机械输出功率}不是简单等于砂动能增长率。

\sourcetest{1.3.4}{48}{56}
\problemtext 质量$M$的人手持质量$m$物体，
以初速$v_0$、仰角$\alpha$从地面向前跳跃，
到弹道最高点时将物体以相对人向后的水平速度$u$抛出。
求抛出物体相对不抛时，此人的水平落地点增加距离。
\workings
跳跃全程竖直初速度$v_0\sin\alpha$，
到最高点历时$t_h=v_0\sin\alpha/g$。
抛出前人和物体的水平速度均为$v_0\cos\alpha$。
水平动量在抛出的短暂时间内守恒；
设抛后人速度$v_M$，物体速度为$v_M-u$，
\[
(M+m)v_0\cos\alpha=Mv_M+m(v_M-u).
\]
解得人新增水平速度
\[
\Delta v_M=v_M-v_0\cos\alpha=\frac{mu}{M+m}.
\]
从最高点到落地仍需$t_h$，
故人与\emph{不抛物体}时落地点的差为
\[
\boxed{\Delta x=\frac{muv_0\sin\alpha}{(M+m)g}}.
\]
\crosscheck 题中的$u$是\emph{相对于人}的抛速。
若改成物体对地的抛速，水平动量方程应相应修改。

\sourcetest{1.3.5}{48}{56}
\problemtext 质量$M$的空心物体A在光滑平面上，
内部水平弹簧劲度系数$k$连接内壁，
质量$m$的物体B紧靠弹簧但未固定。
初始弹簧压缩$\Delta x$，全系统由静止释放。
当B与弹簧分开后，求A的反冲速度大小。
\workings
两物体沿水平方向无外力，
分离后动量大小相等而方向相反，
设A反冲速度大小$u$、B速度大小$v$：
\[
Mu=mv.
\]
无摩擦，弹簧初始势能在恢复原长时化为两物体动能，
\[
\frac12k(\Delta x)^2
=\frac12Mu^2+\frac12mv^2.
\]
将$v=Mu/m$代入求得
\[
\boxed{u=\Delta x\sqrt{\frac{km}{M(M+m)}}}.
\]
\crosscheck 极限$M\to\infty$时$u\to0$；
与动量守恒一致。须把A的整个质量作为反冲体质量，
不能只用内壁弹簧质量。

\subsection{自测题1：改错题}

\sourcetest{1.4.1}{48}{56}
\problemtext 轻弹簧原长时挂一质量为$m$的物体，
从静止释放后向下运动。
某解答认为“合力为零时，重力势能减少等于弹性势能增加”，
由$mg y_0=\tfrac12ky_0^2$求出$y_0=2mg/k$，
但又把$y_0$当成\emph{合力为零}的伸长量。指出并改正。
\workings
以向下伸长量$y$为正，
系统的加速度为$a=g-ky/m$，
所以\emph{合力为零}时
\[
\boxed{y_{\rm eq}=\frac{mg}{k}},
\]
该点是静态平衡点，却不是从原长释放过程的
最大下降点；到这里时速度尚不为零。
机械能守恒（初始$y=0,v=0$）给
\[
\frac12mv^2+\frac12ky^2-mgy=0.
\]
在平衡点代入$y=mg/k$得
\[
\boxed{\frac12mv_{\rm eq}^2=\frac{m^2g^2}{2k}>0}.
\]
而下降的\emph{最低点}满足$v=0$，
由$mg y_{\max}=\frac12ky_{\max}^2$解非零根，
\[
\boxed{y_{\max}=\frac{2mg}{k}}.
\]
\crosscheck 数值$2mg/k$对应\emph{速度为零的最低点}，
并非\emph{合力为零的平衡点}；
原解的错误是将这两个不同事件混为一谈。
% === END guide-selftest01-solutions.tex ===
